% =============================================================================
\section{Discussion and perspectives}
\label{sec:discussion}
% =============================================================================

\subsection{Exploratory research and architectural inductive bias}
\label{subsec:exploratory_inductive_bias}

The modeling architecture was not fixed a priori. Instead, several surrogate formulations were investigated to determine which representations and inductive biases are appropriate for the evolution of high-dimensional beam distributions. The FNO, Transformer, windowed Transformer, and GNO should therefore be interpreted as different hypotheses about the structure of the learned evolution operator rather than as unrelated competing models.

The comparison indicates that the organization of information transfer is central to predictive performance. The dense causal Transformer permits direct interaction with the complete upstream sequence, whereas the windowed Transformer restricts this interaction to a finite causal neighbourhood. The GNO imposes a related but more explicit structure through graph-based message passing and a learned kernel. The improved performance of the windowed Transformer and GNO suggests that unrestricted global mixing is not required to represent the dominant latent dynamics and may introduce weakly constrained couplings during rollout.

This result is physically plausible because beam transport is generated through the composition of lattice transformations and structured collective interactions. Wakefield effects may be nonlocal, but their dependence is directional and constrained rather than arbitrary. A restricted interaction support can therefore provide a useful inductive bias without implying that the complete dynamics are strictly local. Repeated interactions across network depth can still transmit information over longer distances.

The optimum at an intermediate attention window is consistent with this interpretation. A very narrow window may omit relevant multi-element dependencies, while an excessively broad one gradually recovers unrestricted mixing. The results therefore identify the support and geometry of latent information transfer as important modeling choices. They do not establish the universal superiority of one architecture, but show that architectures aligned with meaningful structure in the physical problem can provide more accurate and efficient surrogates.

\subsection{Conditional latent uncertainty}
\label{subsec:discussion_conditional_latent_uncertainty}

Sampling from \(q_{\phi}(z_0\mid X_0,c)\) can be propagated through the complete CVAE--GNO pipeline, producing spatially resolved variability maps and uncertainty estimates for the predicted physical statistics. The inferred variability is structured and concentrated mainly in regions containing significant probability mass, rather than appearing as spatially uniform noise.

The principal result, however, is that the ensemble is under-dispersed relative to the observed prediction error. In the fixed-sample experiment, the mean ensemble-mean discrepancy was approximately \(19.6\) times the mean latent-sampling standard deviation in normalized-mass units. The corresponding ratio in physical-density units was \(55.5\), although this value is affected by inverse-bin-area rescaling. In the independent held-out example, the discrepancy remained approximately \(7.44\) times larger than the sampled spread. Latent sampling therefore explains only a limited fraction of the total surrogate error.

The physical-statistics head shows the same limitation. For one validation example, the predicted energy-deviation spread was only approximately \(11.8\%\) of its reference value, while its latent-sampling variability was negligible in comparison. Since the model weights remain fixed throughout the ensemble, the experiment captures sensitivity to the conditional latent state but excludes parameter uncertainty, finite-data effects, architecture uncertainty, model discrepancy, and out-of-distribution behaviour. The resulting maps may therefore be used as relative sensitivity indicators, but not as calibrated confidence intervals. More complete uncertainty quantification would require additional model variability and empirical coverage tests, for example through independently trained ensembles or post-hoc calibration.

\subsection{Toward accelerator digital twins and optimization}
\label{subsec:digital_twins_optimization}

Fast differentiable surrogates can support tasks beyond direct reproduction of simulation outputs, including sensitivity analysis, parameter scans, model calibration, and optimization. In this sense, the present work contributes to the broader development of accelerator digital twins: reduced models that complement high-fidelity simulation and experimental measurements.

A natural extension is to combine the surrogate with Bayesian optimization. The learned model could provide inexpensive evaluations of beam-quality objectives, while a higher-level optimizer explores machine parameters under operational constraints \cite{roussel_bayesian_2024}. Such a workflow is particularly relevant when direct simulation or machine evaluation is costly. Its reliability would nevertheless depend on improved uncertainty calibration and explicit detection of extrapolative conditions.

\subsection{Limitations}
\label{subsec:discussion_limitations}

Several limitations define the scope of the conclusions. First, the models are validated only within the sampled parameter domain, and generalization outside this region cannot be assumed. Second, the fifteen pairwise projections retain substantial distributional information but do not uniquely determine the complete six-dimensional density. Third, CVAE compression may smooth fine filamentary or weak tail structures. Fourth, the present architectures do not exactly enforce symplecticity, charge conservation, or analytical wakefield constraints.

Finally, the uncertainty ensemble varies only the conditional latent state and does not represent the complete predictive uncertainty. Further work should therefore combine broader and more realistic simulation datasets with physics-constrained architectures, out-of-distribution evaluation, calibrated uncertainty estimates, and comparisons against higher-fidelity collective-effects models. These extensions are necessary before the surrogate can be considered for reliable optimization or operational decision making.